\documentclass[11pt]{article}

\usepackage[margin=0.9in]{geometry}
\usepackage{amsmath,amssymb}
\usepackage{booktabs}
\usepackage{enumitem}
\usepackage{listings}
\usepackage{xcolor}
\usepackage{hyperref}
\usepackage{tikz}

\definecolor{backcolour}{rgb}{0.95,0.95,0.92}
\definecolor{codeblue}{rgb}{0.05,0.15,0.55}
\definecolor{codegreen}{rgb}{0.0,0.45,0.0}

\lstset{
    backgroundcolor=\color{backcolour},
    basicstyle=\ttfamily\small,
    keywordstyle=\color{codeblue}\bfseries,
    commentstyle=\color{codegreen},
    stringstyle=\color{purple},
    breaklines=true,
    frame=single,
    showstringspaces=false,
    tabsize=4
}

\title{CPSC 250 \\ Multiple Inheritance and Mixin Classes}
\author{Edward Brash}
\date{}

\begin{document}

\maketitle

\section{Multiple Inheritance}

Python allows a class to inherit from more than one base class.

This is called \textbf{multiple inheritance}.

For example:

\begin{lstlisting}[language=Python]
class Bat(Winged, Mammal):
    pass
\end{lstlisting}

A bat is both winged and a mammal, so it may make sense for it to inherit from both classes.

\section{Conceptual Example}

Suppose:

\begin{lstlisting}[language=Python]
class Winged:

    def flap_wings(self):
        print("flap flap")


class Mammal:

    def breathe(self):
        print("breathing")

    def give_birth(self):
        print("giving birth")
\end{lstlisting}

Then:

\begin{lstlisting}[language=Python]
class Bat(Winged, Mammal):
    pass
\end{lstlisting}

A \texttt{Bat} object can use methods inherited from both parent classes.

\section{Method Name Conflicts}

A complication arises if both base classes define a method with the same name.

For example, suppose both parent classes define:

\begin{lstlisting}[language=Python]
blink()
\end{lstlisting}

Which version does Python use?

\section{Method Resolution Order}

Python uses the method from the first base class listed, according to its method resolution order.

For example:

\begin{lstlisting}[language=Python]
class Bat(Winged, Mammal):
    pass
\end{lstlisting}

If both \texttt{Winged} and \texttt{Mammal} define \texttt{blink()}, Python will search \texttt{Winged} first.

\section{Mixin Classes}

A mixin is a class designed to add a small piece of behavior to another class.

Mixin classes are not usually intended to stand alone.

They are used to compose behavior.

\section{Transport Example}

Suppose we begin with a base class:

\[
\texttt{TransportMode}
\]

Then we want to create:

\begin{itemize}
    \item a \texttt{Truck},
    \item a \texttt{FlyingCar}.
\end{itemize}

Both are transport modes, but they move differently.

\section{Driving Mixin}

\begin{lstlisting}[language=Python]
class DrivingMixin:

    def drive(self, distance):
        print(f"Driving {distance} miles")
\end{lstlisting}

\section{Flying Mixin}

\begin{lstlisting}[language=Python]
class FlyingMixin:

    def fly(self, distance, altitude):
        print(f"Flying {distance} miles at {altitude} feet")
\end{lstlisting}

\section{Using Mixins}

A truck can drive:

\begin{lstlisting}[language=Python]
class Truck(TransportMode, DrivingMixin):

    def go(self, distance):
        self.drive(distance)
\end{lstlisting}

A flying car can both fly and drive:

\begin{lstlisting}[language=Python]
class FlyingCar(TransportMode, DrivingMixin, FlyingMixin):

    def go(self, distance):
        self.fly(distance / 2, self.altitude)
        self.drive(distance / 2)
\end{lstlisting}

\section{Why Mixins Are Useful}

Mixins allow us to keep behavior modular.

Instead of copying the same \texttt{drive()} method into several classes, we define it once and inherit it where needed.

This supports:

\begin{itemize}
    \item code reuse,
    \item cleaner design,
    \item better encapsulation,
    \item simpler maintenance.
\end{itemize}

\section{Using the Objects}

The end user should be able to write:

\begin{lstlisting}[language=Python]
s = Truck(...)
f = FlyingCar(...)

s.go(100)
f.go(100)
\end{lstlisting}

and have both objects behave appropriately.

This is another example of polymorphism.

\section{Key Takeaways}

\begin{itemize}
    \item Python supports multiple inheritance.
    \item A class can inherit from more than one parent class.
    \item If method names conflict, Python follows method resolution order.
    \item Mixins add focused behavior to classes.
    \item Mixins help avoid duplicate code.
    \item Polymorphism allows different objects to use the same interface.
\end{itemize}

\end{document}
